\section{The Altar of Victory and the Siege of the Basilicas}

On 25~August 383, the usurper Magnus Maximus had Emperor Gratian assassinated at Lyons. The death of Ambrose's orthodox imperial patron left the Western court in Milan in the hands of a twelve-year-old boy, Valentinian~II, and his mother, the Homoian Empress Justina. Ironically, Justina turned first to Ambrose for diplomatic survival: in the winter of 383--384, Ambrose crossed the Alps to Trier on a perilous embassy to Maximus and helped deter an immediate invasion of Italy (\work{Epistle} XX.23; \work{Letter} XXI.20). Yet no sooner had Ambrose helped steady the boy-emperor's throne than two consecutive crises tested the survival of Nicene Christianity in the imperial capital: first the senatorial campaign to restore the Altar of Victory in Rome (384), and then the Milanese court's bid to seize Catholic basilicas for Homoian worship (385--386).

\subsection{Symmachus and the Altar of Victory (384)}

In 382, Emperor Gratian had removed the ancient Altar of Victory from the Roman Senate House, withdrawn public subsidies from the pagan priestly colleges and the Vestal Virgins, and annulled their right to receive legacies in land. After Gratian's murder, pagan senators seized the moment. In the summer of 384, Quintus Aurelius Symmachus---Prefect of the City of Rome, the most admired Latin orator of his day, and a personal acquaintance of Ambrose's family (whom Ambrose had warmly called \latin{parens}, ``kinsman'' or respected elder, in \work{De Excessu} I.32)---submitted his celebrated \work{Relatio} III to Valentinian~II.

Symmachus was no crude anti-Christian fanatic; his memorial deserves to be read in its own noble, elegiac register. Speaking as a Roman magistrate defending constitutional memory, he pleaded that custom bound the civic order together and gave solemnity to senatorial oaths. Summoning the personified figure of aged Rome before the young emperor, he begged: ``Excellent princes, fathers of your country, respect my years to which pious rites have brought me. Let me use the ancestral ceremonies, for I do not repent of them. Let me live after my own fashion, for I am free'' (\work{Memorial of Symmachus} 9). He closed his religious appeal with lines that have echoed through Western philosophy ever since:

\begin{quote}
We ask, then, for peace for the gods of our fathers and of our country. It is just that all worship should be considered as one. We look on the same stars, the sky is common, the same world surrounds us. What difference does it make by what pains each seeks the truth? We cannot attain to so great a secret by one road; but this discussion is rather for persons at ease, we offer now prayers, not conflict. (\work{Memorial of Symmachus} 10)
\end{quote}

Symmachus added two sharp legal arguments: confiscating testamentary lands left to Vestals and priests violated Roman property law, and the severe Italian famine of 383 had followed directly upon this ``sacrilege'' (\work{Memorial} 13--15).

Hearing rumors that Valentinian's consistory was inclined to yield, Ambrose struck immediately with \work{Epistle} XVII before he had even seen Symmachus's text. He denied that Symmachus spoke for the whole Senate, recalling that two years earlier Pope Damasus had sent Ambrose a counter-petition from Christian senators who vowed to boycott the Curia if pagan sacrifices resumed (\work{Epistle} XVII.10). He warned Valentinian~II that a Christian ruler who funded an idol altar made himself a participant in pagan sacrifice, and laid down a blunt warning of episcopal resistance: ``Certainly if anything else is decreed, we bishops cannot contentedly suffer it and take no notice; you indeed may come to the church, but will find either no priest there, or one who will resist you'' (XVII.13).

Once granted a copy of Symmachus's memorial, Ambrose composed his fuller point-by-point reply, \work{Epistle} XVIII. He met Symmachus's eloquence with equally sharp Roman rhetoric, making Rome speak a second time---now as a Christian city unashamed of conversion: ``I do not blush to be converted with the whole world in my old age. It is undoubtedly true that no age is too late to learn. \ldots\ This alone was common to me with the barbarians, that of old I knew not God'' (\work{Epistle} XVIII.7). To Symmachus's pluralist epigram about many roads to the great secret, Ambrose replied with the unmistakable clarity of Christian revelation:

\begin{quote}
By one road, says he, one cannot attain to so great a secret. What you know not, that we know by the voice of God. And what you seek by fancies, we have found out from the very Wisdom and Truth of God. Your ways, therefore, do not agree with ours. You implore peace for your gods from the Emperors, we ask for peace for the Emperors themselves from Christ. (\work{Epistle} XVIII.8)
\end{quote}

He then dismantled Symmachus's financial complaints with courtroom precision. Seven richly endowed Vestal Virgins held purple insignia and public stipends for a fixed term, whereas uncounted Christian virgins lived in voluntary poverty without imperial subsidy (XVIII.11--12). As for inheritance rights, imperial legislation had already barred Christian clergy from receiving legacies from pious widows, yet the Church did not complain: ``I prefer that we should be poorer in money than in grace'' (XVIII.13--14), adding that ``The possessions of the Church are the maintenance of the poor'' (XVIII.16). Finally, he pointed out that the harvest of 384 had already repaid the earth with interest, and that even during the scarcity of 383 Gaul, the Pannonias, Rhaetia Secunda, Liguria, and the Venetias had yielded abundant produce, refuting the claim that the gods had cursed the Empire (XVIII.20--21). Ambrose's letters were read aloud in the imperial consistory; even pagan military officers such as Bauto and Rumoridus yielded to Valentinian's decision to reject Symmachus's petition (\work{Letter} LVII.3).

\subsection{Holy Week 385 and the Basilica Crisis}

No sooner had the Altar of Victory been settled than Empress Justina moved against the Nicene churches of Milan itself. In Holy Week of 385, imperial officials first demanded the Portian Basilica outside the walls, and then the larger New Basilica within the walls, for Homoian Easter worship. Ambrose's contemporary letter to Marcellina (\work{Epistle} XX) captures the hour-by-hour tension with journalistic immediacy.

Refusing the palace counselors---``the temple of God could not be surrendered by a Bishop'' (\work{Epistle} XX.2)---Ambrose kept to his Palm Sunday liturgy, teaching the Creed to baptismal candidates in the baptistery and then beginning to ``celebrate mass'' (\latin{missam facere}, XX.4---the earliest extant literary occurrence of \latin{missa} for the Eucharistic liturgy). During the oblation, word arrived that an angry Catholic street crowd had seized an Arian presbyter named Castulus. Ambrose wept at the altar, prayed that no blood might be shed in the Church's cause, and dispatched priests and deacons to rescue Castulus from the mob (XX.5).

The court retaliated fiercely: merchants of Milan were thrown into chains during the very week debtors were ordinarily released, and fined two hundred pounds of gold within three days; palace officers were ordered off the streets (XX.6--7). When counts and tribunes demanded rapid surrender on the ground that all things belonged to the emperor, Ambrose drew the line between personal property and sacred trust:

\begin{quote}
I answered that if he asked of me what was mine, that is, my land, my money, or whatever of this kind was my own, I would not refuse it, although all that I have belonged to the poor, but that those things which are God's are not subject to the imperial power. ``If my patrimony is required, enter upon it, if my body, I will go at once. Do you wish to cast me into chains, or to give me to death? it will be a pleasure to me. I will not defend myself with throngs of people, nor will I cling to the altars and entreat for my life, but will more gladly be slain myself for the altars.'' (\work{Epistle} XX.8)
\end{quote}

He preached on the appointed Holy Week lessons of Job and Jonah while soldiers ringed the basilica. Remarkably, the blockade collapsed from within: Catholic troops ordered to hold the church entered to pray instead of fighting, and on Holy Thursday the emperor recalled the soldiers and remitted the merchants' fines (XX.13, 20--26). Yet Ambrose warned Marcellina that the storm had not passed. The boy-emperor had cried out to his counts, ``If Ambrose bade you, you would deliver me up to him in chains,'' and the chief chamberlain Calligonus had threatened to take Ambrose's head---drawing Ambrose's stinging reply: ``God grant you to fulfil your threat; for then I shall suffer as bishops do, you will act as do eunuchs'' (XX.27--28).

\subsection{Auxentius of Durostorum and the Constitutional Boundary (386)}

In January 386, the court passed an imperial law granting freedom of assembly to adherents of the 359 Homoian creed of Ariminum and imposing capital penalties on any who opposed them (\work{Letter} XXI.9--11). A rival Homoian bishop, Mercurinus of Scythia---who took the name Auxentius to claim continuity with Ambrose's predecessor---challenged Ambrose to a formal disputation in the imperial consistory before lay judges chosen by both sides, with the unbaptized teenager Valentinian~II presiding as final arbiter.

Ambrose refused the summons in \work{Letter} XXI, citing Valentinian~I's own rescript that ``priests should judge concerning priests'' (XXI.2) and pressing a principle that reverberated across the Middle Ages: ``who is there who can deny that, in a matter of faith,---in a matter I say of faith,---bishops are wont to judge of Christian emperors, not emperors of bishops'' (XXI.4). He would gladly debate Auxentius in church before the Christian people, or at a lawful synod of bishops once the threatening edict was repealed, but he would never hand over judgment on Christ's divinity to a biased palace court (XXI.6, 15--16, 20).

When rumors spread that soldiers were coming to drag him into exile and the faithful barricaded themselves inside the basilica for days and nights singing his new hymns, Ambrose preached his \work{Sermon against Auxentius}. He disavowed violence utterly: ``against weapons, soldiers, Goths, my tears are my weapons, for these are a priest's defence'' (\work{Sermon against Auxentius} 2). He appealed to Naboth's vineyard---refusing to surrender ``the inheritance of Dionysius, who died in exile in the cause of the faith; the inheritance of the Confessor Eustorgius, the inheritance of Myrocles'' (18). He acknowledged Caesar's rightful fiscal domain---``The lands of the Church pay tribute. If the Emperor wants the lands, he has the power to claim them, none of us will interfere'' (33)---and then articulated, in two sentences of crystal clarity, his theology of Church and Empire:

\begin{quote}
We pay to Caesar what is Caesar's, and to God what is God's. Tribute is due to Caesar, we do not deny it. The Church belongs to God, therefore it ought not to be assigned to Caesar. For the temple of God cannot be Caesar's by right. \ldots\ For what is more full of respect than that the Emperor should be called the son of the Church. \ldots\ For the Emperor is within the Church, not above it. (\work{Sermon against Auxentius} 35--36)
\end{quote}

\subsection{The Relics of Saints Gervasius and Protasius (June 386)}

The standoff reached its climax in June 386. Having completed a new basilica (soon called the Basilica Ambrosiana), Ambrose was urged by the congregation to consecrate it with relics of martyrs as he had done at the Roman basilica. ``Certainly I will if I find any relics of martyrs,'' he replied, feeling ``a kind of prophetic ardour'' (\work{Letter} XXII.1). Directing clergy to dig before the chancel screen of the church of Saints Felix and Nabor, he uncovered ``two men of marvellous stature, such as those of ancient days. All the bones were perfect, and there was much blood,'' with heads separated from the bodies; elderly citizens then recalled having once heard the names Protasius and Gervasius and read their titles (\work{Letter} XXII.2, 12).

During the solemn two-day translation through packed crowds from the basilica of Fausta to the Basilica Ambrosiana, Ambrose reported multiple exorcisms and the sudden recovery of sight by a well-known Milanese butcher named Severus, who had lost his sight, been forced to quit his trade, and lived on charity until he touched the hem of the cloth covering the martyrs' bier (\work{Letter} XXII.2, 17). Augustine, residing in Milan at that exact moment, independently corroborated both the discovery and the healing of the blind man twice in his own writings---in \work{Confessions} IX.7.16 (noting that the emperor and whole city knew the man, and that the event restrained Justina's fury) and decades later in \work{The City of God} XXII.8 (``The miracle which was wrought at Milan when I was there, and by which a blind man was restored to sight, could come to the knowledge of many; for not only is the city a large one, but also the emperor was there at the time, and the occurrence was witnessed by an immense concourse of people'').

\witnesslimit{Ambrose's \work{Letter} XXII and Augustine's \work{Confessions} IX.7.16 and \work{City of God} XXII.8 provide rare, independent, contemporary eyewitness attestation that two skeletons were excavated in June 386, identified as Gervasius and Protasius, translated amid massive public devotion, and accompanied by a publicly witnessed healing of the blind butcher Severus that broke the Homoian court's momentum. At the same time, no pre-386 documentary record of the two martyrs survives, Homoian opponents openly mocked the healings at the time (\work{Letter} XXII.16--22), and the later \work{Passio} detailing their martyrdom under Astasius and parentage by St.~Vitalis (reflected in the 1902 \work{Martyrologium Romanum} for 19~June and 28~April) belongs to subsequent hagiographic development rather than fourth-century testimony.}

Ambrose placed the martyrs beneath the right side of the altar in his new basilica, surrendering the vault he had prepared for his own burial: ``He upon the altar, Who suffered for all; they beneath the altar, who were redeemed by His Passion. I had destined this place for myself, for it is fitting that the priest should rest there where he has been wont to offer, but I yield the right hand portion to the sacred victims'' (\work{Letter} XXII.13). Justina's campaign to seize the Milanese churches never recovered.
